%% The first command in your LaTeX source must be the \documentclass command.
\documentclass[]{ceurart}

\sloppy
\usepackage{listings}
\lstset{breaklines=true}

\begin{document}

\copyrightyear{2025}
\copyrightclause{Copyright for this paper by its authors.
  Use permitted under Creative Commons License Attribution 4.0
  International (CC BY 4.0).}

\conference{CLEF 2025}

\title{Job-Skill Matching with Skill Type Classification}

\tnotemark[1]
\tnotetext[1]{Working notes paper submitted to CLEF 2025.}

\author[1]{Hussain Mustansir}[
email=hm07xxx@st.habib.edu.pk,
]
\cormark[1]
\fnmark[1]
\author[1]{Mahrukh Yousuf}[
email=my07xxx@st.habib.edu.pk,
]
\fnmark[1]
\author[1]{Afiya Sohail}[
email=as07xxx@st.habib.edu.pk,
]
\fnmark[1]
\address[1]{Habib University, Karachi, Pakistan}

\cortext[1]{Corresponding author.}
\fntext[1]{These authors contributed equally.}

\begin{abstract}
This literature review addresses the task of Job-Skill Matching with Skill Type Classification.
The task requires, given a job title, retrieving a ranked list of
relevant skills and classifying each as core, complementary, or
transversal. We are provided a training set of at least 5,000 job title--skill pairs, a development
set of 200 job titles with ESCO-normalised skill annotations carrying a core/contextual relevance
label, and a test set of 500 job titles for blind evaluation. We survey
six prior works that address skill extraction, weakly supervised skill matching, occupational
classification, domain-adaptive pretraining, technical NER, and transversal skill detection ---
each contributing methods or findings we draw on in building our approach.
\end{abstract}

\begin{keywords}
  job-skill matching \sep
  skill type classification \sep
  ESCO \sep
  transformer models \sep
  contrastive learning \sep
  named entity recognition \sep
  information retrieval
\end{keywords}

\maketitle

%% ---------------------------------------------------------------
\section{Introduction to the Task}
%% ---------------------------------------------------------------

The task of Job-Skill Matching with Skill Type Classification sits at the intersection of
information retrieval and multi-class labelling. A system must take a raw job title as input,
retrieve a ranked list of skills from the ESCO taxonomy, and simultaneously assign each
retrieved skill one of three labels: \textit{core}, \textit{complementary}, or \textit{transversal}. The training
data offers at least 5,000 job title. The development set introduces a
richer annotation layer of 200 job titles with ESCO-normalised skills marked as either core or
contextual, which forces a model to capture the difference between what a job \textit{always}
demands and what it \textit{sometimes} demands. The challenge lies in combining semantic retrieval 
with skill-type classification under ranking-based evaluation metrics such as nDCG and MAP. 
Because TalentCLEF includes multilingual data, models must operate across languages without relying 
solely on translation. Additionally, skill labels are structured and taxonomy-aligned, requiring 
systems to distinguish between different types of competencies rather than simply detecting skill mentions.

%% ---------------------------------------------------------------
\section{Literature Review}
%% ---------------------------------------------------------------

The work most relevant to this task covers three areas: pulling skills out of job text,
sorting those skills into the right categories, and dealing with skills that are vague or
implied rather than stated outright. Six papers are reviewed below, each approaching one
or more of these problems from a different angle.

Decorte et al.~\cite{Decorte2025ESCO} build a two-stage classifier that detects transversal
skills in job ads using the ESCO taxonomy as the label standard. The first stage assigns a
skill to one of six broad categories---Core, Thinking, Self-management, Social/communication,
Physical, and Life skills. The second stage breaks that category down further into subclasses.
Both stages use 768-dimensional sentence embeddings produced by a BERT-based Sentence
Transformer fed into neural classifiers. One result worth highlighting: a multilingual embedding
model performed within 0.55\% accuracy of an English-only model. That narrow gap is
encouraging for our task, since it means we may not need a separate translation step to
handle non-English job titles. The paper's label structure also maps directly onto the
core/complementary/transversal distinction in TalentCLEF, making it a natural reference
for how we set up the classification layer.

Gugnani and Misra~\cite{Gugnani2020WeaklySupervised} take a different approach to the
extraction problem. Instead of manually labelling training data, they use ESCO itself to
generate labels automatically. The idea is simple: pull all $n$-grams of length one to four
from a job posting, embed them, and compare each embedding against the embeddings of
ESCO skill entries using cosine similarity. Any $n$-gram that clears a set threshold gets
labelled as a skill. They test three ways of building these embeddings on top of both
RoBERTa and the job-domain model JobBERT---encoding the phrase on its own (ISO),
averaging its embedding across every time it appears in the corpus (AOC), and weighting
token embeddings by IDF scores (WSE). WSE comes out on top, hitting a loose F1 of 59.61
with RoBERTa on the Sayfullina benchmark. On the harder SkillSpan benchmark, however,
JobBERT pulls ahead of RoBERTa, which confirms that pre-training on job-specific text
makes a real difference. 

JobBERT-V3, presented by Decorte et al.~\cite{DecorteJobBERTv3}, is trained on 21 million
job titles across English, German, Spanish, and Chinese using GPT-4.1-generated synthetic
translations of a proprietary English dataset. The model is built on multilingual MPNet and
projects its 768-dimensional representations up to 1024 dimensions through an asymmetric
linear layer. It is trained with InfoNCE contrastive loss, where shared ESCO skill annotations
act as the bridge that pulls equivalent job titles from different languages toward the same
point in embedding space. When applied to skill retrieval without any task-specific fine-tuning,
it reaches a MAP of 0.245 on the validation split. 

Bhola et al.~\cite{Bhola2020SkillBERT} train SkillBERT, a BERT model trained from scratch
on 700,000 job requisitions covering 2,997 unique skills, specifically to produce useful
skill embeddings. The downstream task is to classify each skill into one or more of 40
competency groups, making it a multi-label problem. Their feature set combines three types
of signal: cosine similarity scores between a skill's embedding and each group centroid,
cluster membership labels generated by K-means (45 clusters) and spectral clustering
(35 clusters), and the output probabilities from a two-layer neural network classifier
trained on the SkillBERT embeddings directly. The results show a clear jump in XGBoost
F1 when moving from pre-trained BERT (60.84\%) to Word2Vec (66.73\%) to SkillBERT
(90.35\%). The main takeaway is that the quality of the embedding model matters far more than
the cleverness of the features built on top of it. For us, this is a strong case for starting
from a job-domain checkpoint rather than a general BERT, even if training a full SkillBERT
from scratch is not feasible here.

Pooja et al.~\cite{Pooja2024TechnicalNER} focus on extracting technical skills from data
science job postings scraped from Indeed, covering 6,325 vacancies. They compare TF-IDF,
text classification, and Named Entity Recognition (NER) as candidate methods, settling on
NER because it reads context and grammatical structure rather than just counting word
frequencies. They fine-tune three pre-trained models: spaCy, BERT in its RoBERTa form,
and XLNet. BERT gives the best overall result with an F1 of 93.17\%, supported by a
precision of 92.31\% and a recall of 94.19\%. XLNet hits the highest precision at 95.22\%
but recall falls to 69.96\%, which is a meaningful trade-off when missing a skill costs more
than a false alarm. The paper also makes a practical point about preprocessing: dropping
stopwords and applying stemming, steps that help traditional models, actually hurt
transformer-based models because transformers need the full sentence to build accurate
contextual representations. That finding directly shapes the preprocessing choices we make.

Vera et al.~\cite{Vera2025BETO} build an occupational classifier for the Chilean labor market,
mapping job postings to codes in the four-level CIUO-08/ISCO-08 taxonomy. Their dataset
of 4,605 postings, drawn from six online job boards and labeled by domain experts
The model is BETO, a Spanish-language BERT, fine-tuned on this data. At the most detailed (four-digit)
level it reaches a weighted accuracy of 0.69 and an F1 of 0.66, beating both an SVM and an
LSTM baseline. The SVM holds up reasonably at broad levels but falls apart at the fine-grained
ones; the LSTM underperforms across the board. The main structural issue this paper surfaces
is cascading errors: if a prediction is wrong at the top level of a hierarchy, every level
below it is also likely to be wrong. Our three-way classification (core/complementary/transversal)
does not have four levels, but core is implicitly the most selective of the three labels, and
the lesson here is that models which account for that structure tend to do better than
those that treat all labels as equivalent and independent.

Across all six papers, a few consistent takeaways emerge. Domain-adapted encoders
reliably beat general-purpose BERT for both extraction and classification. Cosine similarity
against ESCO embeddings is a more robust retrieval signal than keyword matching. Transversal
skills need their own dedicated treatment and do not respond well to the same NER pipelines
used for technical skills. And benchmark scores should be read with some caution, since
label noise in evaluation sets is a documented problem in this space.

%%
%% Define the bibliography file to be used
\bibliography{bibfile}

%%
%% Appendix
\appendix

\section{Online Resources}

The sources for the ceur-art style are available via
\begin{itemize}
\item \href{https://github.com/yamadharma/ceurart}{GitHub},
\item
  \href{https://www.overleaf.com/latex/templates/template-for-submissions-to-ceur-workshop-proceedings-ceur-ws-dot-org/pkfscdkgkhcq}{Overleaf template}.
\end{itemize}

\end{document}